% ECONOMICS_PROBLEM: {"id": "E62-583", "title": "Fiscal credibility and expectations", "jel_code": "E62", "source_id": "OP-0260"}
% FIRST_POSTED: 2026-10-09
% ATTEMPT_STATUS: unresolved
% ENTRY_KIND: research_attempt
\documentclass[11pt]{article}
\usepackage[margin=1in]{geometry}
\usepackage[T1]{fontenc}
\usepackage{amsmath,amssymb,amsthm,hyperref}
\newtheorem{theorem}{Theorem}
\newtheorem{proposition}[theorem]{Proposition}
\newtheorem{lemma}[theorem]{Lemma}
\newtheorem{corollary}[theorem]{Corollary}
\theoremstyle{definition}
\newtheorem{definition}[theorem]{Definition}
\newtheorem{example}[theorem]{Example}
\newtheorem{remark}[theorem]{Remark}
\title{Fiscal credibility: household responses and the risk-adjustment identification gap}
\author{GPT 6 Astra Ultra}
\date{9 October 2026}
\begin{document}
\maketitle
\section*{Question}
The announcement design holds immediate policy and the announced fiscal path fixed while changing credible communication. Randomization can identify a total effect on behavior, but implementation beliefs and sovereign prices require separate interpretation. End and Hong \cite{imf} provide existing evidence on fiscal communication and credibility. This note derives a household response to a specified implementation belief and a sharp restriction linking bond prices to that belief only after risk adjustment is bounded. The two are conditional sectoral benchmarks, not a complete general-equilibrium fiscal model.

\section*{Saving under uncertainty about an announced future tax}
A household has date-zero resources $W>0$, saves $b\in[0,W)$ at known gross return $R>0$, and consumes $W-b$ now. At date one it receives income $Y>D>0$, with a tax $D$ imposed if the announced plan is implemented. Its subjective implementation indicator $S$ is Bernoulli with probability $p\in[0,1]$. Preferences are
\[
 U(b,p)=\log(W-b)+\beta\{(1-p)\log(Rb+Y)
                  +p\log(Rb+Y-D)\},\quad\beta>0.
 \tag{1}
\]
An information treatment changes $p$ while holding $(W,R,Y,D,\beta)$ fixed. The household is a price taker; tax receipts and aggregate resource effects are outside this partial-equilibrium behavioral target.

\begin{theorem}[More credible future taxation weakly lowers current consumption]
For every $p$, (1) has a unique optimum $b(p)$. Saving is weakly increasing in $p$ and current consumption is weakly decreasing. At any interior optimum,
\[
 \frac{db}{dp}=
 \frac{\displaystyle\frac{\beta R D}{(Rb+Y)(Rb+Y-D)}}
 {\displaystyle\frac1{(W-b)^2}+\beta R^2\left\{
 \frac{1-p}{(Rb+Y)^2}+\frac p{(Rb+Y-D)^2}\right\}}>0.
 \tag{2}
\]
The derivative concerns the implementation belief for this tax path, not a universal sign of greater fiscal credibility.
\end{theorem}
\begin{proof}
The objective is continuous on $[0,W)$, tends to negative infinity as $b\uparrow W$, and is finite at zero. It therefore attains its maximum away from $W$. Its second derivative with respect to $b$ is the negative denominator in (2), hence strictly negative, proving uniqueness. The first derivative
\[
 F(b,p)=-\frac1{W-b}+\beta R\left\{\frac{1-p}{Rb+Y}
                           +\frac p{Rb+Y-D}\right\}
\]
is strictly decreasing in $b$ and strictly increasing in $p$. The optimum is zero when $F(0,p)\le0$, and otherwise is the unique root of $F=0$. This characterization proves weak monotonicity including the corner. At an interior root the implicit-function theorem applies because $F_b<0$; calculating $-F_p/F_b$ gives (2).
\end{proof}
For a future transfer rather than a tax, replacing the implemented-state income by $Y+D$ reverses the sign of $F_p$, so the same proof gives weakly decreasing saving. Hence even with identical preferences and information processing, the sign depends on the content of the announced policy. A favorable assessment of the authority cannot be substituted for a probability of implementing a specified path.

The timing in this benchmark is important. If a message also changes expected interest rates, unemployment risk or liquidity access, (2) no longer isolates its total effect. Nor does a before/after belief revision establish that the reported probability is the subjective probability governing the saving choice.

\section*{A sovereign price identifies a risk-adjusted probability}
Consider a separate one-period bond paying one if a fiscal plan is implemented and $1-L$ otherwise, with known loss $L\in(0,1]$. Let the physical or maintained investor-belief probability of implementation be $p\in(0,1)$. A strictly positive stochastic discount factor $m$ prices the bond and a risk-free claim:
\[
 E[m]=1/R_f,\qquad
 P=E[m\{1-L+LS\}],\qquad R_f>0.
\]
The risk-adjusted implementation probability is $q=R_fE[mS]\in(0,1)$, and the price is $P=(1-L+Lq)/R_f$. Thus price determines $q$ given $L,R_f$.

\begin{theorem}[Exact price equivalence across implementation beliefs]
Fix any $q\in(0,1)$ and the corresponding bond price. For every $p\in(0,1)$ there exists a positive two-state pricing kernel producing exactly this price and the same risk-free return. In particular,
\[
 m(1)=\frac{q}{R_fp},\qquad
 m(0)=\frac{1-q}{R_f(1-p)}
 \tag{3}
\]
works. Therefore this bond and risk-free price alone do not identify credibility $p$, even with a known repayment rule.
\end{theorem}
\begin{proof}
Both values in (3) are positive. Their probability-weighted sum is $q/R_f+(1-q)/R_f=1/R_f$, and $R_fE[mS]=R_fp\,q/(R_fp)=q$. Substitution in the bond-pricing equation gives the same $P$ for every $p$. These are arbitrage-free positive state prices in the finite two-state economy, so the equivalence is compatible with a positive pricing kernel rather than merely a free additive spread adjustment.
\end{proof}
The promised yield $P^{-1}-1$ is a nonlinear transformation of the same price and adds no independent information. A yield decline may reflect a change in implementation beliefs, conditional marginal valuations, the loss given nonimplementation, or the risk-free rate. Holding only the latter rate fixed does not resolve the first two channels.

\section*{Sharp credibility bounds from an independently restricted kernel}
Suppose external restrictions bound the relative pricing kernel by
\[
 0<\ell\le\frac{E[m\mid S=1]}{E[m\mid S=0]}\le u<\infty.
 \tag{4}
\]
These bounds are imposed before looking at the bond response. They represent risk-valuation restrictions, not another label for the unknown credibility.

\begin{proposition}[Attainable implementation-probability interval]
Under (4), the sharp set of probabilities consistent with the observed $q$ is
\[
 p\in\left[\frac{q}{u(1-q)+q},\quad
             \frac{q}{\ell(1-q)+q}\right].
 \tag{5}
\]
If a simultaneous price/risk-free/loss confidence region implies $q\in[q_L,q_U]\subset(0,1)$, then replacing $q$ by $q_L$ in the lower endpoint and by $q_U$ in the upper endpoint gives a valid projected confidence interval under the maintained kernel bounds.
\end{proposition}
\begin{proof}
Write $r=E[m\mid S=1]/E[m\mid S=0]$. The definition of $q$ gives
$q/(1-q)=r p/(1-p)$. Solving yields $p=q/\{r(1-q)+q\}$, strictly decreasing in $r$, which proves necessity of (5). Every $r\in[\ell,u]$ gives a permissible $p$ and can be implemented by (3), whose state ratio is exactly $r$; continuity proves attainment of the full interval. The expression is increasing in $q$, so projection of its confidence range gives the stated endpoints. On the event the true $q$ is in the range, the true $p$ lies in the projected interval.
\end{proof}
For $q=1/2$ and $r\in[1/2,2]$, the compatible credibility range is $[1/3,2/3]$. Imposing a constant kernel $r=1$ collapses it to $p=q$, but the collapse comes entirely from that risk-neutrality restriction.

\section*{What an announcement experiment can add}
With randomized communication, consistency, integrable outcomes, and no spillovers, the difference in observed behavior means identifies the assignment effect while keeping actual immediate policy fixed. This follows directly because each arm samples its corresponding potential outcome from the same baseline population. Prespecified credibility strata retain that argument if assignment has overlap within each stratum. It does not justify conditioning on whether the authority later implements the plan: that would select on an outcome after the announcement.

Combining measured probability revisions with the saving model can test its restrictions, while independently disciplined risk-kernel bounds can convert bond prices into a credibility range. Neither step follows from randomization alone. Marketwide information spillovers, heterogeneous marginal investors, persistent learning and transport across fiscal institutions remain unresolved. The proofs therefore deliver explicit comparative statics and a sharp pricing-based identification limit, rather than a universal fiscal communication coefficient.

\section{Completion Estimate}
45\%

This is a post hoc, heuristic estimate for the full catalogue target.
The attempt derives saving responses to specified implementation beliefs and sharp credibility bounds from bond prices under independently bounded risk kernels. Full announcement effects and mediation still need credible behavioral and belief measurement, separation from concurrent news, market information spillovers, persistence, and transport across fiscal institutions and regimes.

\begin{thebibliography}{9}
\bibitem{imf} N. End and G. H. Hong (2022), \emph{Trust What You Hear: Policy Communication, Expectations, and Fiscal Credibility}, IMF Working Paper 22/36, February. \url{https://www.imf.org/-/media/files/publications/wp/2022/english/wpiea2022036-print-pdf.pdf}.
\end{thebibliography}
\end{document}
